% labshelf: poster
\documentclass[final]{beamer}
\usepackage[T1]{fontenc}
\usepackage[brazil]{babel}
\usepackage[orientation=portrait,size=a0,scale=1.15]{beamerposter}
\usepackage{FiraSans}
\usepackage{newpxtext,newpxmath}
\usepackage{graphicx,booktabs,multicol}
\usepackage[locale=FR]{siunitx}
\usepackage{pgfplots,pgfplotstable}
\usepackage[final]{qrcode}
\usepackage[style=numeric-comp,backend=biber,maxbibnames=3,doi=false,url=false,eprint=false,isbn=false]{biblatex}
\usetikzlibrary{positioning,calc}
\usefonttheme{serif}
% labshelf:style
\definecolor{lsPrimary}{HTML}{0F3D3E}
\definecolor{lsAccent}{HTML}{B4532A}
\definecolor{lsBackground}{HTML}{FBF8F1}
\definecolor{lsText}{HTML}{1C1C1C}
\definecolor{lsMuted}{HTML}{5F6368}
\def\lsScale{1}
\iftutex
  \usepackage{fontspec}
  \setmainfont{TeX Gyre Pagella}
  \setsansfont{Fira Sans}
  \setmonofont{Fira Mono}
  \newfontfamily\lsHeadingFont{Fira Sans}
\else
  \let\lsHeadingFont\relax
\fi
\setbeamercolor{normal text}{fg=lsText,bg=lsBackground}
\setbeamercolor{background canvas}{bg=lsBackground}
\setbeamercolor{structure}{fg=lsPrimary}
\setbeamercolor{frametitle}{fg=lsPrimary,bg=lsBackground}
\setbeamercolor{title}{fg=lsPrimary,bg=lsBackground}
\setbeamercolor{block title}{fg=lsBackground,bg=lsPrimary}
\setbeamercolor{alerted text}{fg=lsAccent}
\setbeamercolor{footline}{fg=lsMuted}
\setbeamerfont{frametitle}{family=\lsHeadingFont}
\setbeamerfont{title}{family=\lsHeadingFont}
% labshelf:style end
% Desenho do pôster sobre o bloco de estilo, só com as cores ls*.
\setbeamertemplate{navigation symbols}{}
\setbeamercolor{block title}{fg=lsPrimary,bg=lsBackground}
\setbeamercolor{block alerted title}{fg=lsBackground,bg=lsPrimary}
\setbeamercolor{block alerted body}{fg=lsText,bg=lsPrimary!6}
\setbeamerfont{block title}{family=\sffamily,series=\bfseries,size=\LARGE}
\setbeamertemplate{block begin}{%
  \vskip1.4cm
  {\usebeamercolor[fg]{block title}\usebeamerfont{block title}\MakeUppercase{\insertblocktitle}\par}
  \vskip0.35cm{\color{lsAccent}\rule{4cm}{4pt}}\vskip0.7cm
  \usebeamercolor[fg]{normal text}}
\setbeamertemplate{block end}{\par}
\setbeamertemplate{block alerted begin}{%
  \vskip0.6cm
  \begin{beamercolorbox}[wd=\linewidth,sep=1.4cm,colsep*=0pt]{block alerted body}
  {\color{lsAccent}\sffamily\bfseries\large\MakeUppercase{\insertblocktitle}\par}\vskip0.5cm}
\setbeamertemplate{block alerted end}{\end{beamercolorbox}\vskip0.6cm}
\setbeamertemplate{itemize item}{\color{lsAccent}$\bullet$}
\setbeamertemplate{caption}{\raggedright\color{lsMuted}\small\insertcaption\par}
\setbeamertemplate{headline}{%
  \begin{tikzpicture}[remember picture,overlay]
    \fill[lsPrimary] (current page.north west) rectangle ([yshift=-15cm]current page.north east);
    \fill[lsAccent] ([yshift=-15cm]current page.north west) rectangle ([yshift=-15.6cm]current page.north east);
    \node[anchor=west] at ([xshift=4cm,yshift=-7.5cm]current page.north west) {\includegraphics[height=9cm]{logo-seminario}};
    \node[anchor=east] at ([xshift=-4cm,yshift=-7.5cm]current page.north east) {\includegraphics[height=8cm]{logo-labshelf}};
    \node[anchor=center,text=lsBackground] at ([yshift=-7.3cm]current page.north) {%
      \parbox{0.66\paperwidth}{\centering
      {\sffamily\bfseries\fontsize{96}{110}\selectfont\inserttitle\par}\vskip1.1cm
      {\sffamily\Large\insertauthor\par}\vskip0.4cm
      {\large\color{lsBackground!75}\insertinstitute\par}}};
  \end{tikzpicture}\vskip15.6cm}
\setbeamertemplate{footline}{%
  \hbox to \paperwidth{\hskip3cm{\color{lsMuted}\small Seminário de Teoria dos Números · dezembro de 2026}\hfill
    {\color{lsMuted}\small código e dados: \texttt{scripts/gerar\_figuras.py}}\hskip3cm}\vskip2cm}
\hypersetup{colorlinks,allcolors=lsAccent}
\pgfplotsset{compat=1.18,
  poster/.style={axis lines=left, axis line style={lsMuted,line width=1.2pt}, tick style={lsMuted},
    tick label style={font=\small,color=lsMuted}, label style={font=\normalsize}, width=\linewidth,
    legend style={draw=none,fill=none,font=\small}, every axis plot/.append style={line width=2.2pt}}}
\renewcommand*{\bibfont}{\footnotesize}

\addbibresource{../../refs.bib}
\graphicspath{{../../figures/}}
\newcommand{\tablesdir}{../../paper/tables}
\input{../../paper/shared/macros}
\input{../../paper/shared/enunciados}

\title{Primos em progressões aritméticas}
\author{Luca Marinho Nasser}
\institute{Seminário de Teoria dos Números \textperiodcentered{} biblioteca LabShelf}

\begin{document}
\begin{frame}[t]

\begin{columns}[T]
  \begin{column}{0.92\paperwidth}
    \begin{alertblock}{A mensagem}
      \begin{columns}[c]
        \begin{column}{0.5\linewidth}
          {\sffamily\fontsize{62}{74}\selectfont\color{lsPrimary}
          Os primos têm \textbf{densidade zero}, mas contêm progressões aritméticas de \textbf{todo comprimento}.\par}
          \vskip1cm
          {\large\color{lsMuted}Green \& Tao~\cite{green2008primes}. Até $10^5$, a contagem real fica a menos de $2\%$ da previsão de Hardy--Littlewood.}
        \end{column}
        \begin{column}{0.46\linewidth}
          \includegraphics[width=\linewidth]{fig-grade210}\par
          {\small\color{lsMuted}Coluna 199 de uma grade de 210: os dez primos $199 + 210n$.}
        \end{column}
      \end{columns}
    \end{alertblock}
  \end{column}
\end{columns}

\begin{columns}[T]
  \begin{column}{0.29\paperwidth}
    \begin{block}{Resumo}
      \input{../../paper/shared/resumo}
    \end{block}
    \begin{block}{Dirichlet}
      \textbf{Teorema.} \enunciadoDirichlet
      \begin{align*}
        \sum_{\substack{p \le x\\ p \equiv a\ (q)}} \log p
          &= \frac{1}{\varphi(q)} \sum_{\chi} \overline{\chi}(a) \sum_{p \le x} \chi(p)\log p \\
          &= \frac{x}{\varphi(q)} + o(x) .
      \end{align*}
      As classes empatam no limite~\cite{dirichlet1837}, mas a classe $3 \bmod 4$ lidera a corrida até $x = 26\,861$ e quase sempre depois~\cite{rubinstein1994}.
      \vskip0.8cm
      \begin{tikzpicture}
        \begin{axis}[poster, height=10cm, xmin=0, xmax=1000000, scaled x ticks=false,
            xtick={0,500000,1000000}, xticklabels={0,$5\cdot10^5$,$10^6$}, xlabel={$x$},
            ylabel={$\pi(x;4,3)-\pi(x;4,1)$}]
          \addplot[lsMuted, line width=1pt] coordinates {(0,0)(1000000,0)};
          \addplot[lsPrimary, line width=1.4pt] table[x=x,y=delta,col sep=comma] {../../paper/data/corrida.csv};
        \end{axis}
      \end{tikzpicture}
    \end{block}
    \begin{block}{Equidistribuição}
      Até $10^7$ as quatro classes módulo $10$ recebem quase exatamente $\pi(x)/4$ primos.
      \vskip0.6cm
      \begin{tikzpicture}
        \begin{axis}[poster, height=8.5cm, ybar, bar width=1.6cm, symbolic x coords={1,3,7,9}, xtick=data,
            xlabel={classe $a$ módulo $10$}, ylabel={desvio de $\pi(x)/4$}, axis x line=bottom,
            nodes near coords, nodes near coords style={font=\small,color=lsMuted}, point meta=rawy,
            every node near coord/.append style={/pgf/number format/fixed,/pgf/number format/precision=0}]
          \addplot[fill=lsPrimary, draw=none] table[x=a,y=desvio,col sep=comma] {../../paper/data/residuos.csv};
        \end{axis}
      \end{tikzpicture}
    \end{block}
  \end{column}

  \begin{column}{0.29\paperwidth}
    \begin{block}{Teoremas}
      \textbf{Szemerédi}~\cite{szemeredi1975}. \enunciadoSzemeredi
      \vskip0.6cm
      \textbf{Green--Tao}~\cite{green2008primes}. \enunciadoGreenTao
      \vskip0.6cm
      \textbf{Szemerédi relativo.} \enunciadoRelativo
    \end{block}
    \begin{block}{Método}
      \input{blocks/metodo}
    \end{block}
    \begin{block}{Razões mínimas}
      Se $p \le k$ e $p \nmid d$, algum termo $a + jd$ é múltiplo de $p$, então $k\# \mid d$ (salvo se $a = p$).
      \vskip0.6cm
      \begin{center}\input{\tablesdir/diferencas}\end{center}
    \end{block}
  \end{column}

  \begin{column}{0.29\paperwidth}
    \begin{block}{Contagem}
      A previsão de Hardy--Littlewood~\cite{hardylittlewood1923} para progressões de $k$ primos até $X$:
      \begin{equation*}
        N_k(X) \sim \mathfrak{S}_k \sum_{a+(k-1)d \le X} \, \prod_{j=0}^{k-1} \frac{1}{\log(a+jd)} .
      \end{equation*}
      \begin{tikzpicture}
        \begin{axis}[poster, height=12cm, xmode=log, ymode=log, xlabel={$X$}, ylabel={progressões até $X$},
            legend pos=north west]
          \addplot[lsPrimary, only marks, mark=*, mark size=7pt] table[x=x,y=pa3,col sep=comma] {../../paper/data/contagem.csv};
          \addplot[lsPrimary!50] table[x=x,y=hl3,col sep=comma] {../../paper/data/contagem.csv};
          \addplot[lsAccent, only marks, mark=square*, mark size=6pt] table[x=x,y=pa4,col sep=comma] {../../paper/data/contagem.csv};
          \addplot[lsAccent!50] table[x=x,y=hl4,col sep=comma] {../../paper/data/contagem.csv};
          \legend{$k=3$, previsão, $k=4$, previsão}
        \end{axis}
      \end{tikzpicture}
      Em $X = 10^5$: $\num{2856331}$ contra $\num{2887579}$ ($k=3$) e $\num{389606}$ contra $\num{396467}$ ($k=4$).
    \end{block}
    \begin{block}{Saiba mais}
      \begin{columns}[c]
        \begin{column}{0.6\linewidth}
          Artigo: \href{https://arxiv.org/abs/math/0404188}{arXiv:math/0404188}\\[0.3cm]
          Exposição: \href{https://doi.org/10.4171/EMSS/6}{Conlon, Fox \& Zhao}~\cite{conlon2014}\\[0.3cm]
          Recorde: 27 primos em progressão (PrimeGrid, 2019).
        \end{column}
        \begin{column}{0.34\linewidth}
          {\color{lsPrimary}\qrcode[height=7cm]{https://arxiv.org/abs/math/0404188}}
        \end{column}
      \end{columns}
    \end{block}
    \begin{block}{Referências}
      \nocite{vandercorput1939,furstenberg1977,gowers2001,green2010linear}
      \printbibliography[heading=none]
    \end{block}
    \begin{block}{Agradecimentos}
      Aos participantes do Seminário de Teoria dos Números pelas discussões, e a quem mantém o PrimeGrid e o TeX Live.
    \end{block}
  \end{column}
\end{columns}

\end{frame}
\end{document}
